\section{Scope and Sharpness of the Hypotheses}
\label{sec:scope-and-sharpness-of-the-hypotheses}

\subsection{The finite-variance assumptions}
\label{sec:finite-variance-assumptions}

We first return to the hypotheses of the finite-variance positive-density
theorem and separate their roles.
The support-aware classification of
\cref{sec:density-zeros-and-the-linear-equivalence-group} shows that
positivity is not merely a regularity convenience: without it,
Hilbert--Schmidt proximity must be paired with the zero-set row condition,
and boundary geometry may collapse the equivalence group to permutations.
The remaining assumptions control the arbitrary-mixing Fourier/Fisher
mechanism.

The finite-variance positive-density theorem is broad within its regime.
It includes centered, variance-normalized Laplace densities despite their
cusp at the origin, as well as variance-normalized densities proportional to
\(e^{-x^4}\).
Symmetry is not needed in any step of the proof.

For examples allowing both asymmetry and heavy tails, take a density
proportional to
\begin{equation}
 (1+y^2)^{-(\tau+1)/2}
                           e^{\beta\arctan y},
 \qquad \tau>2,\quad \beta\in\R,
 \label{eq:asymmetric-heavy-tail-example}
\end{equation}
and center and rescale to variance one.
Its location score before centering and rescaling is
\[
 \frac{(\tau+1)y-\beta}{1+y^2}.
\]
Both this score and its product with \(y\) are bounded.
The density is positive and smooth.
It has a finite second moment.
Centering and rescaling preserve finiteness of the location and scale
Fisher information.
So the standardized density satisfies \eqref{eq:main-assumptions}.
For \(2<\tau\le4\), its fourth moment is infinite.
When \(\beta\ne0\), the two tails have unequal leading constants.
So centering does not make the density even.
The case \(\beta=0\) is a rescaled Student \(t\) family.
These examples show that neither evenness nor a fourth moment is implicit in
the theorem.

The Gamma example used to illustrate the rotational Fourier defect in
\cref{sec:variance-and-affine-normalization} should therefore be read
through the support-aware theorem, not through the positive-density theorem.
Its centered density vanishes on a half-line.
The Fourier defect still gives a local Hellinger obstruction to rotations,
but equivalence is decided by the half-line support classification:
homogeneous linear equivalence permits only coordinate permutations.
The same support theorem classifies all bounded coordinate laws, including
disconnected supports and laws without densities.
\Cref{cor:restricted-gamma-affine} treats the distinct affine Gamma problem,
where support-preserving scales do allow equivalence.

The moment assumptions play a different role from the boundary assumptions.
Support controls which rows are admissible from the viewpoint of null sets;
the second moment controls whether arbitrary \(\ell^2\) rows define the
measurable linear action used throughout \cref{thm:measure-equivalence}.
It also has a specific role beyond differentiating the expectation of the
product of centered complex exponentials.
It makes every \(\ell^2\) row into a convergent random series.
It supplies uniform integrability of squares of linear forms.
It gives the coefficient-subsequence argument with an independent Gaussian
summand in \cref{lem:unit-row-rigidity}.
Thus the first-moment ODE result is stronger in its moment scope than
\cref{thm:measure-equivalence}.

Returning to the finite-variance argument under
\eqref{eq:main-assumptions}, the Hellinger path estimates give \(L^1\)
convergence and uniform integrability of likelihood ratios.
The finite real Fourier sums are bounded in each coordinate.
Their dimension-free bounds are the \(L^2(\mu_n)\) and generator estimates
proved above.
No variance bound for their aggregate under an arbitrary mixed output law is
needed in the tail-symmetry argument.

\subsection{The symmetric stable branch}
\label{sec:symmetric-stable-branch}

The second-moment boundary does not leave the symmetric stable case
unclassified, but it changes both the coefficient domain and the mixing
cost.
Fix \(0<\alpha<2\), let \(\rho_\alpha\) be the standard symmetric
\(\alpha\)-stable law, and put
\begin{equation}
 \widehat\rho_\alpha(t)=e^{-|t|^\alpha},
 \qquad
 \mu_\alpha=\rho_\alpha^{\otimes\N},
 \qquad
 \mu_{\alpha,n}=\rho_\alpha^{\otimes n}.
 \label{eq:standard-symmetric-stable-law}
\end{equation}
Its density \(f_\alpha\) is smooth and strictly positive, and its location
and scale Fisher information are finite, although its second moment is
infinite~\cite{bobkov2014stable}.
Indeed, \(f_\alpha(x)\asymp(1+|x|)^{-1-\alpha}\) and its location score is
\(O((1+|x|)^{-1})\) in the tails, so both the location score and its product
with \(x\) are square integrable against \(f_\alpha\).
For \(\alpha\le1\), the normalization in
\eqref{eq:standard-symmetric-stable-law} specifies a zero location parameter,
not a finite mean.

A stable row series \(\sum_ja_jX_j\) converges almost surely precisely when
\(a\in\ell^\alpha\).
For a real matrix \(T\) with \(\ell^\alpha\) rows, write
\(\Phi_T^{(\alpha)}\) for its stable row-series realization and
\(\nu_T^{(\alpha)}=(\Phi_T^{(\alpha)})_\#\mu_\alpha\).
We call this row action compatibly invertible if there is a matrix \(S\)
with \(\ell^\alpha\) rows such that the two row-series maps are measurable
and
\[
 \Phi_S^{(\alpha)}\circ\Phi_T^{(\alpha)}=I,
 \qquad
 \Phi_T^{(\alpha)}\circ\Phi_S^{(\alpha)}=I
\]
on common conull carriers for the corresponding orbit laws.
This is the stable analogue of the compatible realization supplied by
\cref{lem:measurable-action}; it is an assumption on the action, not a claim
that every bounded \(\ell^2\) operator has stable row sums.

For a matrix \(M=(m_{ij})\), define the anisotropic mixing cost
\begin{equation}
 \mathfrak e_\alpha(M)
 =\sum_i|m_{ii}-1|^2
  +\sum_j\left(\sum_{i\ne j}|m_{ij}|^2\right)^{\alpha/2}.
 \label{eq:stable-mixing-cost}
\end{equation}
The first term measures coordinate scales.
The second first takes the \(\ell^2\)-norm down each off-diagonal input
column and then its \(\alpha\)-th power across source coordinates.
For the actual stable output marginals, set
\begin{equation}
 E_{\alpha,n}(T)
 =-2\log\Aff\bigl(
       \mu_{\alpha,n},(\pi_n)_\#\nu_T^{(\alpha)}\bigr).
 \label{eq:stable-projected-hellinger-energy}
\end{equation}

\begin{theorem}[Symmetric stable coordinates]
\label{thm:symmetric-stable-equivalence}
Let \(0<\alpha<2\) and let \(T\) have a compatibly invertible stable row
action.  Then
\begin{equation}
 \boxed{\begin{aligned}
 \nu_T^{(\alpha)}\sim\mu_\alpha
 &\quad\Longleftrightarrow\quad
 T|_H\in\GL(H)\ \text{and}\
 \mathfrak e_\alpha(Q^{-1}T)<\infty
       \ \text{for some }Q\in\mathcal P\\
 &\quad\Longleftrightarrow\quad
 \sup_nE_{\alpha,n}(T)<\infty,
 \end{aligned}}
 \label{eq:symmetric-stable-equivalence}
\end{equation}
where \(\mathcal P\) is the full signed-permutation group and
\(\mathcal P_n\) is its finite-dimensional counterpart.
If these conditions fail, then
\(\nu_T^{(\alpha)}\perp\mu_\alpha\) and
\(E_{\alpha,n}(T)\uparrow\infty\).
Conversely, if \(T\in\GL(H)\) satisfies the summability condition in
\eqref{eq:symmetric-stable-equivalence}, that condition itself supplies the
stable row action and a compatible row-series inverse.
The exact stabilizer in this class is \(\mathcal P\).
\end{theorem}

\begin{proof}
We organize the proof around the four points at which the finite-variance
argument must be replaced.

\emph{Dichotomy, energy, and the coefficient domain.}
Put \(g_\alpha=\sqrt{f_\alpha}\) and define the translation affinity
\[
 a_\alpha(u)=\int_\R g_\alpha(x)g_\alpha(x-u)\,dx.
\]
Finite positive location Fisher information, strict positivity of the
density, and decay of translations in \(L^2\) give
\begin{equation}
 1-a_\alpha(u)\asymp_\alpha\min\{u^2,1\}.
 \label{eq:stable-translation-affinity}
\end{equation}
Kakutani's theorem therefore identifies \(\ell^2\) as the exact
translation-equivalence space of \(\mu_\alpha\).
If \(\nu_T^{(\alpha)}\sim\mu_\alpha\), translating before and after the
compatible row action shows that \(T\) and its row inverse map this
translation space into itself.
The closed graph theorem then gives \(T|_H\in\GL(H)\).

The compatible two-sided inverse makes every finite family of output rows
linearly independent.
It therefore has a nonsingular input minor, and the remaining stable inputs
contribute an independent convolution factor; hence every finite output
block has a strictly positive density.
The finite-marginal likelihood ratios consequently form the same mean-one
martingale as in
\cref{lem:finite-marginal-likelihood-ratios}.
The product measure is quasi-invariant and ergodic under \(c_{00}\), while
its row-series image has the same properties under \(T(c_{00})\).
The equivalent--singular theorem for quasi-invariant ergodic measures allows
these two translation spaces to differ and therefore applies
\cite{okazaki1985dichotomy}.
It follows exactly as in \eqref{eq:marginal-affinity-limit} that
\begin{equation}
 \nu_T^{(\alpha)}\sim\mu_\alpha
 \quad\Longleftrightarrow\quad
 \sup_nE_{\alpha,n}(T)<\infty,
 \label{eq:stable-energy-dichotomy}
\end{equation}
with singularity and divergent energy on the other branch.

\emph{Sufficiency.}
For a matrix \(K\), put
\[
 F_\alpha(K)=\sum_j\|Ke_j\|_2^\alpha,
 \qquad
 \mathcal C_\alpha=\{K:F_\alpha(K)<\infty\}.
\]
This column class satisfies
\begin{equation}
 F_\alpha(AK)\le\|A\|_{\mathrm{op}}^\alpha F_\alpha(K),
 \qquad
 F_\alpha(KD)\le\|D\|_{\mathrm{op}}^\alpha F_\alpha(K)
 \label{eq:stable-column-class-bounds}
\end{equation}
for bounded \(A\) and bounded diagonal \(D\), and
\begin{equation}
 \|K^*b\|_\alpha^\alpha
 \le F_\alpha(K)\|b\|_2^\alpha.
 \label{eq:stable-row-admissibility-bound}
\end{equation}
Thus \(K\in\mathcal C_\alpha\) has admissible stable rows.

The finite-dimensional estimate that replaces the matrix Fisher upper
bound is
\begin{equation}
 h^2(A_\#\mu_{\alpha,n},\mu_{\alpha,n})
 \le C_\alpha\min\{\mathfrak e_\alpha(A),1\}
 \label{eq:stable-finite-dimensional-upper-bound}
\end{equation}
for every invertible real \(n\times n\) matrix \(A\), with a constant
independent of \(n\).
To prove it, put \(\beta=\alpha/2\), let \(V_i\) be independent positive
\(\beta\)-stable variables with
\(\E e^{-tV_i}=e^{-t^\beta}\), and write
\(X_i=\sqrt{2V_i}Z_i\) with independent standard Gaussians \(Z_i\).
Conditionally on \(V=\diag(V_i)\), Hellinger data processing and the
Gaussian affinity formula bound the left side of
\eqref{eq:stable-finite-dimensional-upper-bound} by a constant times
\[
 \E\min\{\|V^{-1/2}(A-I)V^{1/2}\|_F^2,1\}.
\]
For \(c_j^2=\sum_{i\ne j}|a_{ij}|^2\), the contribution of source column
\(j\) is
\[
 W_j=V_j\sum_{i\ne j}\frac{|a_{ij}|^2}{V_i}.
\]
The positive stable estimates
\[
 \mathbb P(V>r)\le C_\beta r^{-\beta},
 \qquad
 \E[V;V\le r]\le C_\beta r^{1-\beta},
 \qquad
 \E V^{-1}<\infty
\]
give
\[
 \mathbb P(c_j^2V_j>1)
 +\E[W_j;c_j^2V_j\le1]
 \le C_\alpha c_j^\alpha.
\]
Summing over source columns, while retaining the ordinary quadratic bound
on the diagonal, proves \eqref{eq:stable-finite-dimensional-upper-bound}.

Now remove the measure-preserving \(Q\) and write the resulting matrix as
\(M=D+K\), where \(D=\diag(d_i)\), \((d_i-1)\in\ell^2\), and
\(K\in\mathcal C_\alpha\).
The finitely many indices on which the diagonal of \(M\) is not positive
are absorbed into \(K\).
For
\[
 D_n=I+P_n(D-I)P_n,
 \qquad K_n=P_nKP_n,
 \qquad M_n=D_n+K_n,
\]
we have \(M_n\to M\) in operator norm, so the \(M_n\) are eventually
invertible with uniformly bounded inverses.
Write
\[
 M_n^{-1}=D_n^{-1}+J_n,
 \qquad
 J_n=-M_n^{-1}K_nD_n^{-1}.
\]
Then \(F_\alpha(J_n)\) is uniformly bounded, and
\begin{equation}
 \begin{aligned}
 M_n^{-1}M_m-I
 ={}&D_n^{-1}(D_m-D_n)+J_n(D_m-D_n)\\
    &+M_n^{-1}(K_m-K_n).
 \end{aligned}
 \label{eq:stable-relative-approximants}
\end{equation}
The first term has diagonal \(\ell^2\)-norm tending to zero.
The other two tend to zero in \(\mathcal C_\alpha\) by
\eqref{eq:stable-column-class-bounds},
\(\|D_m-D_n\|_{\mathrm{op}}\to0\), and
\(F_\alpha(K_m-K_n)\to0\); their diagonal square cost tends to zero as
well because \(\mathcal C_\alpha\subset\mathcal S_2\).
Thus \eqref{eq:stable-finite-dimensional-upper-bound}, applied to
\(M_n^{-1}M_m\), makes the square-root likelihoods of
\((M_n)_\#\mu_\alpha\) Cauchy in \(L^2(\mu_\alpha)\).
Their limit has norm one, and convergence of each stable row in
\(\ell^\alpha\) identifies its square with
\(d\nu_M^{(\alpha)}/d\mu_\alpha\).
This proves absolute continuity and then equivalence by
\eqref{eq:stable-energy-dichotomy}.
Finally,
\begin{equation}
 (D+K)^{-1}
 =D^{-1}-(D+K)^{-1}KD^{-1}
 \label{eq:stable-inverse-decomposition}
\end{equation}
and \eqref{eq:stable-column-class-bounds} show that the inverse has the same
admissible form.
The continuity of the stable random functional from \(\ell^\alpha\) to
convergence in probability lets the finite-matrix composition identities
pass to the row-series limits, producing compatible measurable versions.
This also proves the last assertion of the theorem.

\emph{The global signed-permutation match.}
Suppose conversely that \(\nu_T^{(\alpha)}\sim\mu_\alpha\).
We transfer the problem to a finite-variance product without applying the
finite-variance theorem directly to the stable law.
Let \(T^{-*}:=(T^{-1})^*\).
Let \(G_\alpha\) be the unitary Fourier transform of \(g_\alpha\), set
\[
 \lambda(d\xi)=|G_\alpha(\xi)|^2\,d\xi,
 \qquad \Lambda_k=(\lambda^{*k})^{\otimes\N}.
\]
The probability \(\lambda\) is symmetric and has variance
\(\|g_\alpha'\|_2^2\in(0,\infty)\).
Moreover, \(G_\alpha*G_\alpha\) is a positive constant multiple of
\(e^{-|t|^\alpha}\); hence
\(|G_\alpha|^2*|G_\alpha|^2>0\) almost everywhere.
The characteristic function of \(\lambda\) is the translation affinity
\(a_\alpha\).
Polynomial stable tails give, for some \(p,\delta>0\),
\[
 c(1+|t|)^{-p}\le a_\alpha(t)
 \le C(1+|t|)^{-\delta},
\]
and \(a_\alpha'\) is bounded.
Choose an integer \(r\) so large that
\[
 \int_\R t^2\bigl(|a_\alpha(t)|^{2r}
                 +|(a_\alpha^r)'(t)|^2\bigr)\,dt<\infty.
\]
Fourier inversion makes the density of \(\lambda^{*r}\) a bounded-variation
density; three such convolutions have finite location Fisher information by
the Fisher smoothing estimate~\cite{bobkov2014stable}.
Let \(b\) be the density of \(\lambda^{*3r}\), put \(q=b*b\), and let
\(s_b,s_q\) be their location scores.
Independent \(U,V\) with density \(b\) satisfy
\[
 x s_q(x)-1=2\E[U s_b(V)\mid U+V=x].
\]
Their finite variance and the finite Fisher information of \(b\) therefore
give finite scale Fisher information for \(q\).
Thus for the even integer \(k=6r\), the standardized density of
\(\lambda^{*k}\) is positive and satisfies
\eqref{eq:main-assumptions}.
It is non-Gaussian because its characteristic function
\(a_\alpha^k\) has a polynomial lower bound, whereas a nondegenerate
Gaussian characteristic function has Gaussian decay.

Let \(p_n=f_\alpha^{\otimes n}\), let \(q_{T,n}\) be the density of the
true first \(n\) stable outputs, and put
\[
 L_n=\frac{q_{T,n}}{p_n},
 \qquad
 \psi_n=\sqrt{q_{T,n}},
 \qquad
 \Omega_n=g_\alpha^{\otimes n}.
\]
Here \(L_n\) is lifted to the product space and
\(\psi_n=\Omega_n\sqrt{L_n}\).
Equivalence gives
\(r_n:=\sqrt{L_n}\to r:=\sqrt L\) in \(L^2(\mu_\alpha)\).

For each \(n\), take \(k\) tensor copies of \(\psi_n\), apply the
finite-dimensional unitary Fourier transform, square the modulus, and push
the resulting probability on \(\R^{kn}\) forward by adding the \(k\)
frequency coordinates belonging to each site.
Call the resulting probability on \(\R^n\) \(\eta_n\), and extend it to
the full sequence space by
\begin{equation}
 \sigma_n=\eta_n\otimes
       (\lambda^{*k})^{\otimes\{i>n\}}.
 \label{eq:stable-fourier-transfer-approximants}
\end{equation}
The finite-dimensional pushforward \(\eta_n\) has a Lebesgue density, and
the density of \(\lambda^{*k}\) is positive almost everywhere; hence
\(\sigma_n\ll\Lambda_k\).

For \(m>n\), pad \(\psi_n\) by \(m-n\) factors of \(g_\alpha\) before
making the Fourier comparison.
Write \(\|\cdot\|_{\mathrm{TV},1}\) for total variation in the
\(L^1\)-norm convention.  Unitarity, the inequality
\(\||u|^2-|v|^2\|_1\le2\|u-v\|_2\) for unit vectors, and contraction of
total variation under pushforward give
\begin{equation}
 \|\sigma_m-\sigma_n\|_{\mathrm{TV},1}
 \le2\|r_m^{\otimes k}-r_n^{\otimes k}\|_{L^2(\mu_\alpha^{\otimes k})}
 \le2k\|r_m-r_n\|_{L^2(\mu_\alpha)}.
 \label{eq:stable-fourier-transfer-tv-bound}
\end{equation}
Thus \(\sigma_n\) converges in total variation to a probability
\(\sigma\ll\Lambda_k\).

Let \(\nu^{(n)}=L_n\mu_\alpha\); this measure has the true first \(n\)
output marginal and the original independent tail.
For every finitely supported \(h\), Fourier autocorrelation and
\(r_n\to r\) give
\begin{align*}
 \widehat\sigma(h)
 &=\lim_n\Aff\bigl(\nu^{(n)},(x\mapsto x+h)_\#\nu^{(n)}\bigr)^k\\
 &=\Aff\bigl(\nu_T^{(\alpha)},
        (x\mapsto x+h)_\#\nu_T^{(\alpha)}\bigr)^k\\
 &=\prod_i a_\alpha((T^{-1}h)_i)^k.
\end{align*}
The last expression is the characteristic function of
\((T^{-*})_\#\Lambda_k\), so finite-dimensional uniqueness identifies
\(\sigma=(T^{-*})_\#\Lambda_k\).
Consequently
\begin{equation}
 (T^{-*})_\#\Lambda_k\ll\Lambda_k.
 \label{eq:stable-fourier-transfer}
\end{equation}
The equivalent--singular alternative in \cref{thm:measure-equivalence}
upgrades this absolute continuity to equivalence.  Applying that theorem to
the standardized \(\lambda^{*k}\) product then gives a signed permutation
\(Q\) such that
\(T^{-*}-Q\in\mathcal S_2\).
Since
\[
 T-Q=-T(T^{-*}-Q)^*Q,
\]
we obtain \(T-Q\in\mathcal S_2\).
After removing \(Q\), write
\begin{equation}
 Q^{-1}T=I+H,
 \qquad H\in\mathcal S_2.
 \label{eq:stable-hilbert-schmidt-match}
\end{equation}

\emph{The fractional column necessity.}
It remains to strengthen \eqref{eq:stable-hilbert-schmidt-match} from a
quadratic conclusion to the second sum in
\eqref{eq:stable-mixing-cost}.
The local input is the dimension-free rare-source estimate
valid for every dimension and every real square matrix \(B\): there are
\(\varepsilon_\alpha,c_\alpha>0\) such that
\begin{equation}
 \diag B=0,\quad F_\alpha(B)\le\varepsilon_\alpha
 \quad\Longrightarrow\quad
 h^2((I+B)_\#\mu_{\alpha,n},\mu_{\alpha,n})
 \ge c_\alpha F_\alpha(B).
 \label{eq:stable-rare-source-bound}
\end{equation}
Here is the rare-source construction.
Write \(F=F_\alpha(B)\),
\(c_j=\|Be_j\|_2\), and
\(u_j=Be_j/c_j\) when \(c_j\ne0\).
Then \(\|u_j\|_2=1\), \((u_j)_j=0\), and the stable scale contributed to
output row \(i\) satisfies
\begin{equation}
 r_i:=\sum_j|b_{ij}|^\alpha\le F.
 \label{eq:stable-row-scale-bound}
\end{equation}

First, for any finite family of vectors \(v_j\) in a real Hilbert space,
the stable random sum \(S=\sum_jv_jX_j\) obeys
\begin{equation}
 \mathbb P(\|S\|_2>t)
 \le C_\alpha t^{-\alpha}\sum_j\|v_j\|_2^\alpha.
 \label{eq:stable-hilbert-random-sum-tail}
\end{equation}
Indeed, exclude the events
\(\|v_j\|_2|X_j|>t\), truncate on their complement, and apply Chebyshev's
inequality to the centered truncated sum.
The stable tail and truncated second moment give
\eqref{eq:stable-hilbert-random-sum-tail}.
In particular, there is a deterministic
\(\eta_\alpha(s)\downarrow0\) as \(s\downarrow0\) such that
\begin{equation}
 \E\min\left\{\left\|\sum_jv_jX_j\right\|_2,1\right\}
 \le\eta_\alpha\left(\sum_j\|v_j\|_2^\alpha\right).
 \label{eq:stable-hilbert-random-sum-truncation}
\end{equation}

Second, one bounded statistic detects a unit displacement in every
Euclidean direction, uniformly in the dimension.
Take \(\psi(x)=\tanh x\).
There are a smooth odd increasing function \(\chi\), bounded in modulus by
one, and
\(\kappa>0\), depending only on \(\rho_\alpha\), such that for every unit
vector \(u\),
\begin{equation}
 Z_u(x)=\sum_i u_i\psi(x_i),
 \qquad
 R_u(t)=\E\chi(Z_u(X+tu)),
 \qquad
 R_u(t)\ge\kappa\quad(1\le t\le2).
 \label{eq:stable-uniform-receiver}
\end{equation}
To verify this, note that \(Z_u(X)\) has mean zero and a variance independent
of \(u\).
Choose a common interval carrying most of its mass and then choose \(\chi\)
whose derivative has a positive lower bound on the enlarged interval.
Since
\[
 \inf_{|z|\le2}\E\psi'(X+z)>0,
\]
differentiation gives a dimension-free positive lower bound for
\[
 R_u'(t)=\E\left[
   \chi'(Z_u(X+tu))\sum_i u_i^2\psi'(X_i+tu_i)
                 \right]
\]
when \(|t|\le2\), after discarding the common small exceptional event.
Oddness then proves \eqref{eq:stable-uniform-receiver}.

Third, choose a nonzero smooth odd function \(w\), supported where
\(1\le|t|\le2\), with \(0\le w\le1\) on the positive half-line.
Put
\[
 Z_j(x)=\sum_{i\ne j}(u_j)_i\psi(x_i),
 \qquad
 [z]_{[-1,1]}=\max\{-1,\min\{z,1\}\},
\]
and define the bounded test
\begin{equation}
 \Phi_B(x)=
 \left[\sum_{j:c_j\ne0}w(c_jx_j)\chi(Z_j(x))\right]_{[-1,1]}.
 \label{eq:stable-rare-source-test}
\end{equation}
For all sufficiently small \(c_j\), stable tail asymptotics give
\begin{equation}
 \mathbb P(1\le|c_jX_j|\le2)\le C_\alpha c_j^\alpha,
 \qquad
 \E|w(c_jX_j)|\ge c_\alpha c_j^\alpha.
 \label{eq:stable-rare-gate-probability}
\end{equation}
The gates depend on independent input coordinates.
If \(K\) is the number of open gates, then
\(\E K(K-1)\le C_\alpha F^2\).
Before the outer truncation, each summand in
\eqref{eq:stable-rare-source-test} has mean zero because its gate variable
is independent of its receiver variables.
The truncation changes the sum only when at least two gates are open, so
\begin{equation}
 |\E\Phi_B(X)|\le C_\alpha F^2.
 \label{eq:stable-rare-source-input-mean}
\end{equation}

It remains to control the same statistic at \(Y=(I+B)X\).
For fixed \(0<\delta<1/4\),
\eqref{eq:stable-row-scale-bound} and the stable tail bound imply
\begin{equation}
 \mathbb P\{\exists i:\ |c_i(BX)_i|>\delta\}
 \le C_\alpha\delta^{-\alpha}\sum_i c_i^\alpha r_i
 \le C_\alpha\delta^{-\alpha}F^2.
 \label{eq:stable-output-gate-exception}
\end{equation}
On the complementary event, Lipschitz continuity of \(w\) and
\eqref{eq:stable-rare-gate-probability} show that replacing every output
gate \(w(c_jY_j)\) by its input gate \(w(c_jX_j)\) changes the expectation
by at most
\begin{equation}
 C_\alpha\delta F+C_\alpha\delta^{-\alpha}F^2.
 \label{eq:stable-output-gate-error}
\end{equation}
After this replacement, removing the outer truncation costs at most another
\(C_\alpha F^2\).

Finally, set \(B^{(j)}=B(I-e_j\otimes e_j)\).
The random vector \(B^{(j)}X\) is independent of the gate variable \(X_j\),
and \eqref{eq:stable-hilbert-random-sum-truncation} shows that replacing
\(Z_j(Y)\) by \(Z_j(X+c_jX_ju_j)\) costs, after summing over \(j\), at most
\begin{equation}
 C_\alpha F\eta_\alpha(F).
 \label{eq:stable-receiver-interference-error}
\end{equation}
Because \((u_j)_j=0\), conditioning on \(X_j\) and using
\eqref{eq:stable-uniform-receiver,eq:stable-rare-gate-probability} gives
\begin{equation}
 \sum_j\E\bigl[
   w(c_jX_j)\chi(Z_j(X+c_jX_ju_j))\bigr]
 \ge c_\alpha F.
 \label{eq:stable-rare-source-main-response}
\end{equation}
Choose \(\delta\) and then \(F\) sufficiently small.
Equations
\eqref{eq:stable-rare-source-input-mean,eq:stable-output-gate-error,
eq:stable-receiver-interference-error,eq:stable-rare-source-main-response}
give
\[
 \E\Phi_B(Y)-\E\Phi_B(X)\ge c_\alpha F.
\]
The test can be nonzero only when a gate is open.
The coordinates of \(Y\) have stable scales bounded by
\((1+F)^{1/\alpha}\), so
\[
 \E\Phi_B(X)^2+\E\Phi_B(Y)^2\le C_\alpha F.
\]
The Hellinger Cauchy--Schwarz inequality now gives
\(h^2((I+B)_\#\mu_{\alpha,n},\mu_{\alpha,n})\ge c_\alpha F\),
which is \eqref{eq:stable-rare-source-bound}.

We use this estimate through tail sign symmetries, which avoids any
triangularity assumption.
Let \(B\) have \(\ell^\alpha\) rows, zero diagonal,
\(\|B\|_{\mathrm{op}}\le1/10\), and belong to \(\mathcal S_2\).  Suppose
the row action of \(I+B\) and its row-series inverse are compatible and
\((I+B)_\#\mu_\alpha\sim\mu_\alpha\).
For a finite coordinate set \(J\), let \(S_J\) change the signs in \(J\)
and set
\[
 \Delta_J=S_JBS_J-B,
 \qquad
 C_J=(I+B)^{-1}S_J(I+B)S_J.
\]
The matrix \(\Delta_J\) contains precisely the entries crossing the cut
between \(J\) and its complement.
Because \(J\) is finite, these entries form a sum of finitely many
\(\ell^2\) columns and finitely many \(\ell^\alpha\) rows; hence
\(\Delta_J\in\mathcal C_\alpha\).
The finite-dimensional bound
\eqref{eq:stable-rare-source-bound} extends to such infinite matrices by
applying it to \(P_n\Delta_JP_n\) and using the Hellinger convergence from
the already proved sufficiency argument.

Put
\[
 K_J=(I+B)^{-1}\Delta_J,
 \qquad
 D_J=I+\diag K_J,
 \qquad
 \widetilde B_J=D_J^{-1}(I+K_J)-I.
\]
Then \(C_J=I+K_J\) and \(\diag\widetilde B_J=0\).
Since \(\|(I+B)^{-1}-I\|_{\mathrm{op}}\le1/9\), orthogonal projection
away from the diagonal entry in each column gives
\[
 \frac89\|\Delta_Je_j\|_2
 \le\|(I-e_j\otimes e_j)K_Je_j\|_2
 \le\frac{10}{9}\|\Delta_Je_j\|_2.
\]
For small \(F_\alpha(\Delta_J)\), the diagonal entries of \(D_J\) lie in
\([1/2,3/2]\), and consequently
\begin{equation}
 F_\alpha(\widetilde B_J)\asymp_\alpha F_\alpha(\Delta_J),
 \qquad
 \|\diag K_J\|_2
 \le\frac19F_\alpha(\Delta_J)^{1/\alpha}.
 \label{eq:stable-sign-cut-comparison}
\end{equation}
The rare-source lower bound for \(\widetilde B_J\), the diagonal scale upper
bound for \(D_J\), and the Hellinger triangle inequality now give
\begin{equation}
 F_\alpha(\Delta_J)\le\varepsilon_*
 \quad\Longrightarrow\quad
 h((C_J)_\#\mu_\alpha,\mu_\alpha)
 \ge c_*F_\alpha(\Delta_J)^{1/2}.
 \label{eq:stable-sign-cut-lower-bound}
\end{equation}
Indeed, the two terms have respective orders
\(F_\alpha(\Delta_J)^{1/2}\) and
\(F_\alpha(\Delta_J)^{1/\alpha}\), and the latter is smaller for small
cost because \(\alpha<2\).

On the other hand, finite-prefix approximation of the likelihood ratio
gives the uniform tail symmetry
\begin{equation}
 \sup_{\substack{J\subset\{N+1,N+2,\ldots\}\\ |J|<\infty}}
 h\bigl((I+B)_\#\mu_\alpha,
        (S_J)_\#(I+B)_\#\mu_\alpha\bigr)
 \longrightarrow0.
 \label{eq:stable-tail-sign-symmetry}
\end{equation}
Hellinger invariance and \((S_J)_\#\mu_\alpha=\mu_\alpha\) also give the
connecting identity
\begin{equation}
 h\bigl((I+B)_\#\mu_\alpha,
        (S_J)_\#(I+B)_\#\mu_\alpha\bigr)
 =h(\mu_\alpha,(C_J)_\#\mu_\alpha).
 \label{eq:stable-sign-cut-distance-identity}
\end{equation}
If \(F_\alpha(B)=\infty\), deleting finitely many columns changes the cost
by a finite amount, and so does deleting finitely many rows because those
rows lie in \(\ell^\alpha\).
Hence every tail still has infinite column cost.
Let \(E\) be a finite coordinate set in such a tail, and form
\(J\subset E\) by independently including each coordinate with probability
one half.
Column by column, concavity of \(t^{\alpha/2}\) gives
\begin{equation}
 \E_JF_\alpha(P_E\Delta_JP_E)
 \ge2^{\alpha-1}F_\alpha(P_EBP_E).
 \label{eq:stable-random-sign-cut}
\end{equation}
As \(E\) increases through finite subsets of the tail, the right side is
unbounded, so some finite cuts have arbitrarily large cost.
Equivalence also makes every sufficiently remote one-coordinate output
marginal close to \(\rho_\alpha\), so
\[
 r_i:=\sum_{j\ne i}|b_{ij}|^\alpha\longrightarrow0.
\]
Together with \(B\in\mathcal S_2\), which gives
\(c_i:=\|Be_i\|_2\to0\), this yields
\[
 F_\alpha(\Delta_{\{i\}})
 =2^\alpha(c_i^\alpha+r_i)\longrightarrow0.
\]
Set \(\gamma=\max\{1,\alpha\}\) and
\(\mathcal N(A)=F_\alpha(A)^{1/\gamma}\).
For \(\alpha\le1\), subadditivity of the \(\alpha\)-power makes
\(\mathcal N\) subadditive; for \(\alpha\ge1\), it is the
\(\ell^\alpha(\ell^2)\) norm and satisfies the triangle inequality.
Moreover, when \(i\notin J\),
\[
 \Delta_{J\cup\{i\}}-\Delta_J
 =S_J\Delta_{\{i\}}S_J,
\]
and sign multiplication preserves \(\mathcal N\).
Thus the \(\mathcal N\)-increment from adding one sufficiently remote
coordinate tends uniformly to zero.
Fix \(0<\eta<\varepsilon_*/2\).
Starting with the empty cut and adding the elements of a finite cut supplied
by \eqref{eq:stable-random-sign-cut}, stop the first time the corresponding
\(\mathcal N\)-threshold is crossed.
For every sufficiently remote tail this produces a finite \(J\) with
\(\eta\le F_\alpha(\Delta_J)\le2\eta\).
Equations
\eqref{eq:stable-sign-cut-lower-bound,eq:stable-tail-sign-symmetry,
eq:stable-sign-cut-distance-identity} then contradict one another.
Hence \(F_\alpha(B)<\infty\).

Finally, put \(R=I+H\) in
\eqref{eq:stable-hilbert-schmidt-match} and let
\(R_n=I+P_nHP_n\).
For a sufficiently large \(n\), the equivalent transformation
\(A=R_n^{-1}R\) is arbitrarily close to the identity in operator norm.
Indeed, \(R_n\) changes only finitely many coordinates, so
\((R_n)_\#\mu_\alpha\sim\mu_\alpha\); compatible composition with
\(R_\#\mu_\alpha\sim\mu_\alpha\) gives
\(A_\#\mu_\alpha\sim\mu_\alpha\).
Its diagonal \(D=\diag(a_{ii})\) is positive, boundedly invertible, and
\(D-I\in\mathcal S_2\); the one-coordinate scale criterion therefore gives
\(D_\#\mu_\alpha\sim\mu_\alpha\).
Consequently
\[
 D^{-1}A=I+B,
 \qquad
 \diag B=0,\quad B\in\mathcal S_2,\quad
 \|B\|_{\mathrm{op}}\le1/10,\quad
 (I+B)_\#\mu_\alpha\sim\mu_\alpha.
\]
Finite-coordinate and bounded diagonal composition preserve the
\(\ell^\alpha\) row condition and the compatible row action, so this
normalized \(B\) also satisfies the hypotheses of the sign-cut argument.
The sign-cut argument gives \(B\in\mathcal C_\alpha\).
Writing \(R_n=I+F_n\), where \(F_n\) is supported on a finite coordinate
block, we have
\[
 R=R_nD(I+B)
  =D+\bigl(F_nD+DB+F_nDB\bigr).
\]
The term in parentheses belongs to \(\mathcal C_\alpha\) by
\eqref{eq:stable-column-class-bounds} and finite-coordinate support.
Restoring the initial signed permutation therefore gives
\[
 \sum_j\left(\sum_{i\ne j}|h_{ij}|^2\right)^{\alpha/2}<\infty,
\]
while \(H\in\mathcal S_2\) already gives
\(\sum_i|h_{ii}|^2<\infty\).
This is precisely \(\mathfrak e_\alpha(Q^{-1}T)<\infty\), and completes
the necessity.

If \(\nu_T^{(\alpha)}=\mu_\alpha\), comparison of the joint characteristic
functions gives
\(\|T^*u\|_\alpha=\|u\|_\alpha\) for every finitely supported \(u\).
The same statement for the compatible inverse makes \(T^*\) a surjective
isometry of the stable coefficient space.
For \(0<\alpha<2\), the equality case of
\[
 2\|u\|_\alpha^\alpha+2\|v\|_\alpha^\alpha
 -\|u+v\|_\alpha^\alpha-\|u-v\|_\alpha^\alpha\ge0
\]
shows that disjointly supported vectors have disjointly supported images.
The images of the coordinate vectors therefore have pairwise disjoint
supports; surjectivity makes each support a singleton, and norm preservation
makes its nonzero entry \(\pm1\).
Thus \(T\in\mathcal P\), while every member of \(\mathcal P\) plainly
preserves \(\mu_\alpha\).  This proves the stabilizer assertion and completes
the proof.
\end{proof}

\subsection{Anisotropic Hellinger geometry and sharpness}
\label{sec:stable-hellinger-geometry}

The two summands in \eqref{eq:stable-mixing-cost} have genuinely different
homogeneities.
For a finite matrix \(A\), put
\[
 \mathcal E_{\alpha,n}(A)
 =\min_{Q\in\mathcal P_n}\mathfrak e_\alpha(Q^{-1}A).
\]
If all singular values of \(A\) lie in a fixed interval
\([a,b]\subset(0,\infty)\), then constants depending only on
\(\alpha,a,b\), and not on the dimension, satisfy
\begin{equation}
 c_{\alpha,a,b}\min\{\mathcal E_{\alpha,n}(A),1\}
 \le h^2(A_\#\mu_{\alpha,n},\mu_{\alpha,n})
 \le C_{\alpha,a,b}\min\{\mathcal E_{\alpha,n}(A),1\}.
 \label{eq:stable-hellinger-coercivity}
\end{equation}
The upper bound is \eqref{eq:stable-finite-dimensional-upper-bound}, after
removing the minimizing signed permutation.
For the even convolution power \(k\) constructed in the proof of
\cref{thm:symmetric-stable-equivalence}, the finite-dimensional form of the
Fourier transfer is the quantitative inequality
\begin{equation}
 h^2\bigl((A^{-*})_\#(\lambda^{*k})^{\otimes n},
                 (\lambda^{*k})^{\otimes n}\bigr)
 \le k\,h^2(A_\#\mu_{\alpha,n},\mu_{\alpha,n}).
 \label{eq:stable-finite-fourier-transfer}
\end{equation}
where \(A^{-*}=(A^{-1})^*\).
Indeed, the unitary Fourier transform intertwines the half-density actions
of \(A\) and \(A^{-*}\); taking moduli is an \(L^2\) contraction.
For \(k\) tensor copies, squared Hellinger distance grows by at most the
factor \(k\), and coordinatewise frequency addition is a data-processing
map.
Apply the finite-variance coercivity theorem to the standardized
\(\lambda^{*k}\) product.
When the singular values of \(A\) lie in \([a,b]\), those of \(A^{-*}\)
lie in \([b^{-1},a^{-1}]\), and
\[
 A-Q=-A(A^{-*}-Q)^*Q
\]
shows that distance from the signed-permutation group changes by at most a
factor depending on \(a,b\).
Thus \eqref{eq:stable-finite-fourier-transfer} gives
\begin{equation}
 h^2(A_\#\mu_{\alpha,n},\mu_{\alpha,n})
 \ge c_{\alpha,a,b}
     \min\left\{\min_{Q\in\mathcal P_n}\|A-Q\|_F^2,1\right\}.
 \label{eq:stable-quadratic-lower-bound}
\end{equation}

Near a signed permutation, remove that permutation and write
\[
 d=\sum_i|a_{ii}-1|^2,
 \qquad
 F=\sum_j\left(\sum_{i\ne j}|a_{ij}|^2\right)^{\alpha/2}.
\]
For \(d+F\) small, the identity is the nearest signed permutation, so
\eqref{eq:stable-quadratic-lower-bound} controls \(d\).
If \(D=\diag(a_{ii})\) and \(D^{-1}A=I+B\), then
\(\diag B=0\) and \(F_\alpha(B)\asymp_\alpha F\).
The rare-source bound and the diagonal scale upper bound give
\[
 h(A_\#\mu_{\alpha,n},\mu_{\alpha,n})
 \ge c_\alpha\sqrt F-C_\alpha\sqrt d.
\]
If \(d\) dominates \(F\), use
\eqref{eq:stable-quadratic-lower-bound}; otherwise absorb the second term in
the last display.
This proves the local lower bound in
\eqref{eq:stable-hellinger-coercivity}.
If uniform separation failed away from that neighborhood, one could take a
block-diagonal sum of counterexamples whose squared Hellinger distances are
summable.
Kakutani's theorem would make the resulting infinite stable products
equivalent.
The Hilbert--Schmidt part of \cref{thm:symmetric-stable-equivalence} forces
its global matching permutation to preserve every sufficiently late block;
otherwise each crossed block contributes a fixed positive square error.
The restrictions of that permutation to the late blocks would then make
their anisotropic costs summable, a contradiction.
This proves the global lower bound.

The different local powers are already visible in two dimensions.
For the shear \(T_t(x,y)=(x+ty,y)\), conditioning on \(y\) and applying
Plancherel gives the exact formula
\begin{equation}
 h^2((T_t)_\#\mu_{\alpha,2},\mu_{\alpha,2})
 =2\int_\R(1-e^{-|t\xi|^\alpha})
                  |G_\alpha(\xi)|^2\,d\xi.
 \label{eq:stable-shear-hellinger-formula}
\end{equation}
Consequently
\[
 h^2((T_t)_\#\mu_{\alpha,2},\mu_{\alpha,2})
 \sim2|t|^\alpha\int_\R|\xi|^\alpha
                         |G_\alpha(\xi)|^2\,d\xi,
 \qquad t\to0,
\]
whereas a diagonal scale perturbation has squared Hellinger cost of order
\(t^2\).
Thus the stable geometry is not generated by one finite matrix Fisher
quadratic form.

In particular, Hilbert--Schmidt proximity is necessary but no longer
sufficient.
On independent two-coordinate blocks, let
\[
 T_a(x_k,y_k)_{k\ge1}=(x_k+a_ky_k,y_k)_{k\ge1}.
\]
Then \(T_a-I\in\mathcal S_2\) exactly when \(a\in\ell^2\), while
\cref{thm:symmetric-stable-equivalence} gives
\((T_a)_\#\mu_\alpha\sim\mu_\alpha\) exactly when
\(a\in\ell^\alpha\).
Any \(a\in\ell^2\setminus\ell^\alpha\) therefore gives a
Hilbert--Schmidt perturbation with singular image law, even though the
stable coordinate density is everywhere positive and has finite location
and scale Fisher information.
Conversely, when one source is mixed into many outputs, the criterion
charges the \(\ell^2\)-norm of that whole column only once.
An entrywise \(\ell^\alpha\) condition would therefore be too strong.
The correct stable criterion is an anisotropic, coordinate-sensitive cost,
not a Schatten-class replacement for the finite-variance theorem.
